% ============================================================
%  G. Brandes, "Dyret i Mennesket"
%  Af Dagens Krønike, Maj–Juni 1890, pp. 494–551
%  (Reprinted in: Samlede Skrifter, vol. VII)
%  TRANSCRIPTION (Danish)
% ============================================================
\documentclass[12pt,a4paper]{article}
\usepackage[utf8]{inputenc}
\usepackage[T1]{fontenc}
\usepackage[danish]{babel}
\usepackage{microtype}
\usepackage{setspace}
\usepackage[margin=1.2in]{geometry}
\usepackage{fancyhdr}
\usepackage{hyperref}

\hypersetup{
  pdftitle={Dyret i Mennesket},
  pdfauthor={G. Brandes},
  hidelinks
}

\pagestyle{fancy}
\fancyhf{}
\rhead{Brandes, \emph{Dyret i Mennesket} (1890)}
\cfoot{\thepage}

\setstretch{1.3}

\title{\textbf{Dyret i Mennesket}}
\author{G.\ Brandes}
\date{1890}

\begin{document}
\maketitle
\thispagestyle{fancy}

\begin{center}
  \small (Bourget: \textit{Disciplen}. Maupassant: \textit{Den unyttige Skønhed}.
  Tolstoj: \textit{Kreutzersonaten}. Zola: \textit{Menneskedyret}.)
\end{center}

\bigskip

\noindent Det er paafaldende, i hvor høj Grad den sidste Tids Digtning
hos dens mest fremragende europæiske Repræsentanter er sysselsat med
Forestillingen om en Dobbelthed i vort Væsen.

Hos en enkelt fransk Forfatter Paul Bourget fremtræder denne Dobbelthed
i Menneskevæsenet videnskabeligt begrundet. Bourget er nemlig ogsaa som
Digter den sjælegranskende Kritiker, han begyndte med at være. I
Forklaringen af sine Personers Driftsliv og Følelsesliv gaar han aldrig
rent kunstnerisk til Værks; han peger stadigt med Ransagerens Sonde.

Hvis man overskuer hvad han digterisk har frembragt lige fra hans første
Novelle \textit{L'irréparable} til hans sidste Roman \textit{Le disciple},
er der derfor visse paafaldende fælles Træk. Alt drejer sig i Grunden om
Jegets Fordobling. Hans første Novelle er Historien om en ejendommelig ung
Pige, hvis Livs kritiske Begivenhed er en Overrumpling, hun i Nattens
Ensomhed har været Genstand for fra en samvittighedsløs Mands Side, med
hvem hun har indladt sig paa et uskyldigt Koketteri. Fortællingen begynder
meget betegnende med en Samtale mellem Forfatteren og en Professor i
Sjælelære, der har skrevet et Værk om Opløsning af Idéforbindelserne.

Da nu det, vi kalder vort Jeg, denne Forestilling om noget Indre og
Blivende, der er forenet med et Legem, er en Forestilling, som opstaar
ved Sammenknytning af Ideer, saa er det klart, at hvad Professoren i
Virkeligheden har studeret er Formerne for Jegets Opløsning, for dets
Dobbelthed, som denne optræder under Viljens og Sindets talrige Sygdomme.
I den sidste Roman, Bourget har udgivet, møder vi paany Professoren i
Sjælelære og paany det samme Emne. Det Arbejde, som den unge Forbryder
Robert Greslou, der er anklaget for Mord, har indsendt til en berømt
Professor for at erfare hans Mening, er et Studie over Jegets Dobbelthed.

I Novellen udvikles først, hvori Jegets simpleste Fordobling bestaar.
Jeget er dels bevidst, dels ubevidst. Deraf forklares i Almindelighed
Fejltagelser og Fejlgreb i Livet. Der er skjult i os en Skabning, som vi
ikke kender og om hvilken vi aldrig véd, om den ikke netop er det modsatte
af det Væsen, vi tror at være. Herpaa beror de forunderlige Omslag i
Befindende og Handlemaade, som vi oplever eller iagttager. En arbejder
f.\ Eks.\ mod et Maal, af hvilket han indbilder sig at hans Lykke
afhænger, og naar det er naaet, opdager han, at han har miskendt sit
Følelseslivs hemmelige og sande Fordringer. En gaar og ser ud som et roligt
Menneske, og handler saaledes; pludselig rejser Billedet af hans Ulykke
sig for ham, en stadig Sorg eller stadig Fare, han søger at glemme, men
som kan overvælde ham og give hans Liv en forandret Retning.

I Romanen leverer Hovedpersonen en fuldstændig Undersøgelse af sit eget
Væsen, idet han tilbagefører enhver Tilstand i sit Sind til et
videnskabeligt Skema: Det Nedarvede, Omgivelsernes Indflydelse,
Omplantning af Personen osv. Han bestemmer sin herskende Egenskab som Evne
og Trang til Selvfordobling: Der er altid hos ham ligesom to tydeligt
adskilte Personligheder, en som kommer og gaar, handler og føler, og en
der som Tilskuer iagttager den andens Livsførelse. Han er ikke i Stand til
at angive, hvilken af disse to Personer der er hans sande Jeg. Han forsøger
at udlede denne Dobbelthed hos sig af en Raceblanding, af Nedstamning fra
et Grænsefolk. Han har fra Barn af havt sin Glæde af Forstillelse og
Usandhed, ikke for at prale, men simpelthen for i sin Fantasi og i
Omgivelsernes Bevidsthed at være en Anden. Han har den største Fornøjelse
af at udtale Meninger, som er ganske modsatte dem, han anser for de sande,
og det af den samme besynderlige Grund. Hans Lyst er den at spille en
Rolle ved Siden af sin sande Natur. Den hele Løgnagtighed i hans Holdning,
hans Glæde ved at bedrage, forføre, overliste og tilrane sig en ung
ædelttænkende Kvinde, for hvem han ingen Ømhed nærer. Alt er en Følge af
den for ham uovervindelige Dobbelthed i hans Menneskevæsens oprindelige
Anlæg.

Som man ser, er den Tvedelthed, der sysselsætter den videnskabeligt
dannede Bourget, en saadan, som kun ad filosofisk Vej tilfulde forstaas.
Den svarer ganske vist nogenlunde til den Spaltning af vort Væsen, som man
i gamle Dage fandt imellem det saakaldt Legemlige og det saakaldt
Sjælelige, men den falder paa ingen Maade sammen med den. Kun forsaavidt
opfattes den gammeldags, som Bourget i sine seneste Skrifter synes at ville
forherlige den religiøse Andagt som det eneste Middel imod den.

\section*{I}

Andre af Tidsalderens betydelige Forfattere, der ligeledes dvæler ved
Forestillingen om en Dobbelthed af vort Væsen, har givet dette Tvedelte
et Udtryk, som nærmest hører Middelalderen til. De udtaler sig ikke
sjældent om det Dyriske i Mennesket og om det Menneskelige i Mennesket
som om to forskellige Verdener, der er sammenkoblede i os. De viser sig
optagne af dette Forhold omtrent som Teologerne af Forholdet mellem den
guddommelige og den menneskelige Natur i Gudmennesket. For ikke ganske
faa af de samtidige Digtere er den menneskelige Skabning en Art Centaur,
halvt et højere Væsen, halvt et Dyr, der, naar den fortaber sig i
højtflyvende Tanker, vækkes op af Drømmen ved Lyden af sine egne
Hovslag.

De, der grubler over Dyret i Mennesket, vil enten blot skildre det,
udmale det i dets Vælde og Afskyelighed, eller de vil bekæmpe det,
stundom endog for at udrydde det, uden at det dog lykkes dem nogensinde
skarpt at paavise den Skillelinje, hvor det Dyriske hører op og det
Menneskelige begynder. Der er jo kun altfor mange Virksomheder og
Drifter, som Mennesket har tilfælles med de lavere Kreaturer.

Tankevækkende er nu den principielle Uenighed imellem vor Tidsalders
betydeligste Forfattere angaaende det Spørgsmaal, om det saakaldte
Dyriske er en uovervunden Levning af den oprindelige Naturtilstand, det
vilde Tidehvervs og Barbariets Dage, eller om det, som vi kender det, er
en Kulturfrembringelse, fostret af Overkulturen, saa vi alene kan bekæmpe
det ved, som det hedder, at vende tilbage til Naturen, det vil sige til
simplere, landligere Tilstande.

En Digter som Maupassant eller Zola er nærmest af den første Mening; for
dem er Dyret den oprindelige Vildskab. En Digter som Strindberg eller som
Tolstoj hylder den anden Synsmaade; for dem er Dyret vor falske
Civilisations Unatur.

Tolstojs nye, i Rusland forbudte og meget læste Fortælling
\textit{Kreutzersonaten}, er et rent moraliserende Arbejde, forsaavidt
beslægtet med hans opdragende Legender og Historier for den russiske
Menigmand. Kun at Bogen henvender sig til en Læsekres, der hører til de
mest oplyste i Verden.

Her som i \textit{Gengangere} og \textit{En Handske} er det Køns-Renhed og
-Urenhed, hvorom Alt drejer sig; men der er i den store Russers Synsmaade
en vildere, lidenskabeligere Følgestrenghed end hos Nordmændene. Ibsen
havde i \textit{Gengangere} malt forfærdelige Følger af Mandens Letsind
for Hjemmet og den næste Slægt; Bjørnson havde ladet en ung Kvinde stille
Fordringer til Mandens ubetingede Afholdenhed før Ægteskabet; det syntes
ikke muligt at naa højere eller videre i sine Krav paa dette Felt. Men
Tolstoj løber som ægte Russer Linen ud --- dette er det Særmærke, hvorved
jeg i min Bog \textit{Indtryk fra Rusland} har forsøgt at male russisk
Ejendommelighed. Han genoptager Fordringen til Mandens Renhed før
Ægteskabet, men han forøger den med en Fordring til saakaldt Renhed i
Ægteskabet, der er ny.

Ganske vist taler Tolstoj i sin ny Fortælling ikke i eget Navn. Men dels
gennem en Samtale i en Jernbanevogn, i hvilken de Meninger, der bliver de
sejrrige, synes at maatte ligge nær op til Digterens egne, dels gennem
Hovedpersonens lange Redegørelse for sit Liv, som optager Størstedelen af
Værket, spores Tolstojs Samhu med den Livsanskuelse, der kommer til Orde.
Det er ikke sandsynligt at han skulde ville hævde en væsenlig
Uoverensstemmelse mellem Posdnysjevs Grundsynsmaade og sin egen. I ethvert
Tilfælde har han intet gjort for at bibringe Læseren det Indtryk, at en
saadan Uoverensstemmelse finder Sted. Alt hvad man kan sige er, at han i
sit Skrift \textit{Hvori bestaar min Tro?} vel har bevæget sig ad de samme
Baner som hans Hovedperson, men ligefuldt ikke er gaaet saa vidt.

Tolstoj hævder med Lidenskab den gamle Kristendoms Opfattelse af den
jomfruelige Stand som den egenlig høje og værdifulde. Han har kraftige Ord
om det Taabelige i en Samfundsaand, der kaster mangen ung elskværdig Pige
i en uren og fordærvet Mands Arme blot for at hun ikke skal lide den
formentlige Skam at forblive Jomfru, hvad i Tolstojs Øjne er en Ære.
Allerede dette er radikalt nok følt, om Sagen end ikke er tænkt igennem.

Ægteskab eller Ikke-Ægteskab, Vielse eller Ikke-Vielse er for Tolstoj
ganske ligegyldigt; paa Formen lægger han ingen Vægt. Men den første
Kvinde, en Mand har bundet sig til; den første Mand, en Kvinde har givet
sin Tro, er han og hun knyttet til for Livstid, uopløseligt. De er
sammenknyttede i et Pligtforhold. Om Glæden derved spørges ikke; der
ligger ingen Vægt paa den.

Den moderne Anskuelse, at det er Kærlighed, der begrunder Ægteskabet og
giver det Værd, synes Tolstoj en Letfærdighed, omtrent som den forekom
Fortidens Mennesker, ja i Grunden endnu Hegel med hans saglige Sindelag.
Tolstoj synes at være naaet til den efter adskillige mørktskuende
Menneskers Mening rigtige Overbevisning, at hvad man tidligere kaldte
Ægteskab og mest endnu kalder saadan, kun er muligt, hvor Alt er bygget
paa Myndighed og Lydighed, hvor Forældrene vælger Datteren en Mand,
skaffer Sønnen en Hustru, naar Tidspunktet dertil er kommet, og hvor
Konen skælver for den Husbond, som det er hendes Pligt at elske. Dette er
jo nemlig Ægteskabets klassiske Form: intet Valg, ingen Skilsmisse; én
Herre, én Vilje, i Opsætsighedstilfælde Prygl; megen Hustugt. Vil man
ikke dette, synes han at mene, saa naaer man lidt efter lidt til at gøre
Kærligheden eneraadig, gøre Kønsforbindelsen til en Privatsag og
indskrænke Samfundsvirksomheden til efter Evne at sikre Børnenes Vel.

Den mest konservative men ikke fyldestgørende Fortolkning af
\textit{Kreutzersonaten} er den, at Tolstoj har villet sige: «Som Mændene
nu er, gøres de ved deres Ungdomslevned uskikkede til ægteskabeligt
Samliv.» Det er utvivlsomt, at han, ganske som Bjørnson, er meget opfyldt
af det Slette og Nedværdigende i det Liv, som Mændene efter Sigende i
Reglen fører, forinden de indgaar Ægteskab. Tolstoj har endog i sit
Skriftemaal selv givet Bekendelser af ret uhyggelig Natur, der om de
endogsaa overdriver, viser at han af egen Erfaring kender, hvad han
dømmer haardest:

\begin{quote}
«Det er mig umuligt at mindes hine Aar uden at der overfalder mig en
Følelse af Rædsel, Afsky og Smerte. Der gives ikke den Last, som jeg jo
den Gang kunde have hengivet mig til, ikke den Forbrydelse, som jeg jo
vilde været i Stand til at begaa. Løgn, Tyveri, alskens Udsvævelser,
Vold og Drab --- det er noget som jeg altsammen har gjort mig skyldig i
\ldots\ Ude paa Godset, hvor jeg boede, forødte jeg i Sus og Dus og
Kortspil Alt hvad mine Livegne havde erhvervet til mig ved deres Arbejd
--- tillige pinte og straffede jeg dem ved enhver Lejlighed, ofrede dem
for mine Udsvævelser, bedrog dem, solgte dem osv.»
\end{quote}

Man sammenligne hermed i \textit{Kreutzersonaten} Posdnysjevs Skildring
af, hvorledes han har «væltet sig i Udsvævelsernes Skarn». Kærnen heri er
som i Bjørnsons Fremstillinger af beslægtet Art: Mændene, alle Mænd af de
bedrestillede Klasser idetmindste, lever i deres Ungdom et rentud svinsk
Liv, raat, kærlighedsløst, med Kvinder iflæng.

Det Visse af Sagen er, at den Art Skildringer og Udtalelser i Grunden
alene kan have en levende Interesse for dem, hvis Udvikling ikke er løbet
disse Digteres forbi. Et ikke ganske ringe Antal af mandlige Læseres
Indtryk vil være dette: Vi føler os ikke ramte. Mangen En, der ikke er
mistænkt for at være nogen Hykler, vil i sit stille Sind sige: Jeg er
træt af denne Tale. Det er jo Tale for den aandelige Almue. Hvor
fortræffeligt slige Bøger end kan være skrevne, i dette Punkt er det
Almuebøger, som kun paany indprenter En, at al virkelig Kunst henvender
sig til de mest Udviklede i Tiden.

Vi ser i de forskellige Lande beslægtede literære Forekomster af denne
Art. Dèr opstaar den yngre Dumas, der, efter hvad han selv fortæller, i
sin Ungdom ingenlunde har levet som en Helgen «med mindre man vilde tænke
paa den hellige Augustinus i dennes første Livsafsnit»; han ordner sit
Levned og præker paa Tærsklen til Alderdommen den strengeste Kønsmoral
--- spækker den kun af gammel Vane med dristige Talemaader og slibrige
Hentydninger. Dèr opstaar Bjørnson, som uden at have Dumas' parisiske
Ungdomserfaringer har udviklet sig til en ligesaa haardnakket Pligtlærer
og som har det af en ægte Nordmand, at han er mere frisk og oprigtig end
egenlig fin. Han forkynder os Mændenes kønslige Forfald. Han har, synes
det, set saare megen Raahed omkring sig, og kan han formindske den, er
det fortjenstfuldt. Om den formindskes ved en rundrejsende Moral-Agitation
som den, han i sin Tid slog ind paa, er et aabent Spørgsmaal. Sikkert er
det, at han har forringet sine store opdragende Romaners kunstneriske
Værdi ved den Sjælesorg, der behersker dem.

Nu kommer Tolstoj. Han har en russisk «Junkers» civile og militære
Erfaringer bag sig og han drager Slutninger ud fra disse overfor Mænd,
der aldrig har været «Junkere», aldrig ført Junkerliv, maaske aldrig
udvist Raahed i Forholdet til det andet Køn. Men paa saadanne Læsere
gør nødvendigvis han som de andre Sædelighedsforkyndere i deres Kamp mod
Samfundsmoralen Indtryk af Puritanere, der angriber Papister. Den udviklede
Læser interesserer sig ikke for indre teologiske Fejder. Munkekævl! som
Ulrich von Hutten skrev om Luthers første Angreb paa Paven.

Der gives sikkert endnu i vore Dage (eller allerede i vore Dage) Mænd,
der paa afgørende Punkter føler som Grækere fra det gamle Hellas's bedste
Tid, Mænd, for hvem Tvedelingen af det Dyriske og det Menneskelige ikke
er til, Mænd, der aldrig har væltet sig i Pølen, ingen ond Samvittighed
har og ingen Omvendelse har behov. Disse Digtere kommer nu til dem og
behandler dem som Drankere af Fag: For Himlens Skyld, tilraaber de
Læserne, ikke en Draabe Vin over jere Læber! Ser jer om! I lever i en
Verden af Drukkenbolte og Delirister. Se, dèr ligger En i Rendestenen
døddrukken indtil Umenneskelighed; dèr ligger En skælvende med dansende
Muskelfibre, dødssyg i en Celle af Madratser. For Himlens Skyld, tøm
intet Glas Vin! Og imens modnes Druer under Solens Straaler; de hænger
der, mørke og lyse, og Dionysos, Vinens og Tragediens Gud, staar med det
grønne Løv om Panden og rækker den Tørstige og Trætte Drikken, der er
Menneskets Trøst. Disse Moralister gør Bakkus selv til en
Deliriumskandidat. Og de glemmer, at deres Tale ikke sjældent vender sig
til Mænd, der føler som Grækere, dyrker Vinguden, drikker Vinen blandet
med Vand, nyder den lette Rus og skyr Drukkenskaben som en Væmmelse og
en Pest.

Dog Tolstoj nøjes ikke med det Bjørnsonske Standpunkt. Han gaar meget
grundigere tilværks. Op af hans Bog slaar en mægtig Følelsestroesiver
som i Flammer; og hans følgestrenge Tanke gaar saa meget dybere end
Bjørnsons Forkyndelse, som Oldkristendom er dybere end Udviklingslærens
Seminarie-Visdom.

Tolstoj forarges over Dyret i Mennesket, over Menneskedyret. Han forarges
over Ægtefællernes Samliv i Ægteskabet. Det lidenskabelige Samliv mellem
Nygifte nedbryder efter hans Forestilling al inderligere Velvilje, al
virkelig Agtelse mellem Mand og Kvinde. Det er Skyld i alt Ondt, der
senere dukker op dem imellem. Det afstedkommer Selvforagt og Foragt for
den anden Part. Det vækker i Mellemtiden imellem den legemlige
Tiltræknings Tidsafsnit det dybe Had, der efter hans Opfattelse i Reglen
kommer til Udbrud mellem Ægtefæller i enhver Trætte selv om den ringeste
Genstand, og det er desuden Skyld i den halvt vanvittige Skinsyge, der
efter Bogens Forudsætning altid fortærer en af Parterne eller begge og
gør deres huslige Liv til et evigt Helved.

Det er morsomt i vore Dage at træffe en saadan oldkristelig Følelsesmaade.
Intet Spor af den sunde, antike og moderne Livsopfattelse, der i Samlivet
mellem to Elskende, hvor lidenskabeligt det end monne være, ser noget
Rent og Smukt. Intet Spor af den Naturdyrkelse, der i Oldtiden adlede
Sanselivet; intet Spor heller af den uendelige Ømhed, der hos moderne
Mennesker adler det og forandrer den raa Fornøjelse til Glæde eller Lykke.

Udgangspunktet er Bibelordet, at den, der blot ser paa en Kvinde med
Attraa, han synder med hende i sit Hjerte. Og det Nye er, at for Tolstoj
gælder dette Ord ogsaa i Ægteskabet. Ægteskabet brydes ved Ægtefællernes
Attraa til hinanden. I det fuldkomne Forhold er Kvinden for Manden
Søster, ikke andet. Den klare og rigtige Følge vilde den være, som de
religiøse Særtroende Skoptserne i Rusland drager, nemlig
Selvlemlæstelsen. Tolstoj skyer øjensynligt den. Men til Gengæld ønsker
han, saavidt man kan forstaa, det sanselige Element enten ganske dødet
eller dog trængt ned til det mindst mulige. Ikke alene enhver Blotstillen
af Legemets Former ved nedringede eller tætsluttende Dragter, ikke blot
Dans er ham en Gru som Sanseophidselse; men han frygter Kunsten
overhovedet, særlig Musiken, mindre som Appel til Sanserne end som
Følelsesvækker. Det er da ogsaa efter et Musikværk at Bogen har Navn.
Tankegangen er den, at enhver betydelig Komposition henriver mægtigt, ja
som ved Besættelse Tilhøreren til den bevægede Stemning, i hvilken
Komponisten befandt sig, da han frembragte den; Tilhørerens og
Komponistens Sjæle flyder sammen. Og saaledes sammensmelter da ogsaa
Musiken Tilhørerne, undertiden to og to, i denne samme Stemning, skiller
dem fra Omgivelserne og forener dem. Den kan forsaavidt gøre
Koblergerning.

\section*{II}

Følelsen af selve Kønsforholdet som en Nedværdigelse er efter Europas saa
langvarige Opdragelse i kristelig-aandig Retning ingen sjælden Foreteelse.
Man møder den imidlertid i vore Dage lige saa udviklet hos Personligheder
af ganske hedensk Type som hos Aandsdyrkere og Kristne. Den moderne
Udvikling har ført forfinede Mennesker af de forskelligste Overbevisninger
lige langt bort fra den græske Enhedsfølelse i dens Frihed og Sundhed.
Jeg mindes fra min tidlige Ungdom en lang filosofisk Samtale med en ung
fransk Ven, der drejede sig om dette samme Punkt, Kønsforeningen, som
oprørte ham. Han harmedes over, at Verdensbygmesteren ikke havde fundet
paa en finere, ædlere Form for Kønnenes Sammensmeltning. Et Udraab af ham
ligger mig endnu i Øret: \textit{Pourquoi cette insulte!} Hvortil denne
Forhaanelse, dette Slag i Ansigtet!

Det er denne samme Tanke, der paa flere Steder er varieret i Maupassants
sidste Bog \textit{L'inutile beauté}. I \textit{Un cas de divorce} kommer
Helten overfor sin unge Hustru tilsidst saa vidt, at han ikke kan røre
ved hende med sin Haand eller sine Læber uden at føle en uovervindelig
Væmmelse, ikke just ved hende, men en mere omfattende og mere foragtrig
Væmmelse ved selve det erotiske Favntag, der for alle forfinede Væsener
staar som en Ting, man maa skamme sig over og skjule, ja som man kun
omtaler halvhøjt, med Rødmen.

Og da i den Novelle, hvorefter Samlingen har Navn, Talen falder paa en
smuk, fin Dame, hvis Skønhed har lidt under talrige Barselsenge, og den
Bemærkning falder: Hvad er dertil at sige, det er Naturens Gang! udbryder
den unge Mand, hvem Forfatteren bruger som Talerør: «Naturen! jeg siger
dig, at Naturen er vor Fjende, at vi altid maa kæmpe imod Naturen; den
fører os stedse tilbage til Dyret.» Og nu udvikler han, at alt hvad der
er renligt, smukt, udsøgt, idealt paa Jorden, det er ikke Guds men
Menneskets Værk:

\begin{quote}
«Det er os, som ved at besynge Skabningen, fortolke den, beundre den som
Digtere, forædle den som Kunstnere, forklare den som Videnskabsmænd, har
indført en Smule Ynde og Skønhed, noget Tiltrækkende og
Hemmelighedsfuldt i den. Tænk paa Forplantningen! Hvad kan tænkes mere
uædelt og modbydeligt end denne smudsige og latterlige Akt, over hvilken
alle finere Sjæle til Dagenes Ende vil føle sig oprørte. Dersom Organerne
nødvendigvis skulde tjene et dobbelt Formaal, hvorfor har Skaberens da
ikke til Forplantningen, denne ædleste og vigtigste af alle menneskelige
Virksomheder, kunnet vælge Organer, som var mindre urene, mindre
besudlede? Det er som havde han i en Art ondskabsfuld Cynisme villet
lægge Manden uoverstigelige Hindringer i Vejen for nogensinde at kunne
forskønne eller forædle sit Forhold til Kvinden. Manden har ikke desmindre
opfundet Elskoven, hvad ikke er saa daarligt som Svar til en snu og
drillende Guddom.»
\end{quote}

Her hos Maupassant er altsaa i dette Forhold Naturen opfattet som Fjenden.
Hos Tolstoj er den det nærmeste Maal, dog et Maal hvorudover han atter
stræber ud; et bag os liggende Ideal, hvortil der bør vendes tilbage, kun
at han kender et endnu højere Ideal end den, nemlig Fornegtelsen af den.
Hos ham er Kampen mod Kønslivet kun en Anvendelse af hans almindelige
Synsmaade, ifølge hvilken Kulturen er af det Onde: Penge, Pragtbygninger,
Forfinelse, Overflod, alt andet Liv end Bonde- og Haandværkerliv af det
Onde. Han havde allerede forlængst forsmaaet og fordømt Videnskaben. Nu
er Turen kommen til Kunsten. Han havde i Penge aldrig villet se fortættet
Arbejde, kun fortættet Voldsmagt. Nu ser han i det kønslige Samliv kun en
dyrisk og nedværdigende Tilstand. Han har grublet over dette Tragiske og
Komiske, at Mennesket efter sit Væsen er et videre udviklet Dyr, saa det
kun delvis kan bringe sin Oprindelse i Glemme, aldrig ganske udslette
Sporene af den. Og han ser vor Kultur som den Magt, der giver alt det
Dyriske Overtaget. Man føler i den med Aarene hos ham stigende Afsky for
det blot Naturlige, det Altformenneskelige, baade russisk Religiøsitet og
det Oldingeblik paa disse Mysterier, der har gjort sig gældende ogsaa i
Norden. Afskyen for Letfærdigheden, for det forhøjede Driftsliv, for
Lidenskaben og Lasten har hos ham antaget den endelige Form af Væmmelse
ved Dyret i Mennesket.

Gaar man nu gennem Bogen til dens Forfatter og fra ham til det store
Publikum, der fortaber sig i Beundring for ham, synes følgende Bemærkning
mig at ligge nær. Det er forbavsende som Menneskeheden endnu den Dag
idag bliver slaaet af Ærefrygt for Pligtlærere. Det ligger Flertallet af
os fra den bibelske Opdragelse saaledes i Blodet, at vi uvilkaarligt
mener: Her, hvor Moralen prækes, her foregaar noget særlig Alvorligt og
Højtidsfuldt.

Der havde Russerne en fin, elskværdig, godhjertet Forfatter som
Turgéniev, en Verdensmand, der fordringsløst skrev det ene digteriske
Mesterværk efter det andet. Han indbildte sig ikke at omdanne
Menneskeslægten derved. Han var en stilfærdig Kunstner med enkelte
Svagheder, med store Dyder, en Gentleman.

Saa kommer efter ham en Skribent som Tolstoj. Han er i sin Ungdom en ret
udsvævende Godsejer og Officer, bliver i sin modnere Alder en stor
formanende Digter, med Aarene en stedse mere lidenskabelig Pligtlærer,
tilsidst en Johannes den Døber. Han former sig sin egen Religion, han
præker, han kalder og revser, han gaar i Kamelhaarsskjorte eller rettere
i Musjik-Kofte. Han lader Skægget vokse og skiller Haaret midt i Panden
som en russisk Bonde, han lever paa sit Landgods, fordybet i Betragtning
og Haandgerning; snart ligner han en Bondeprofet; og da hans daglige
Syssels Alvor omsider har præget sig dybt i hans Miner, tillader han en
tilrejsende Fotograf at fastholde dette Udtryk paa den forberedte Plade.
Han gaar bag Ploven med sit hvide Øg for den gammeldags Træplov, han
drømmer sig tilbage i Naturtilstanden, langt bort fra Kulturen, fra hans
Herresæde, der ligger lige i Nærheden, og han tillader Riepin at male
ham, saaledes tilbagesat i Naturtilstanden.

Forskellen er da tydelig nok mellem ham og den oprindelige Johannes den
Døber. Denne havde som bekendt intet Herresæde, saa han havde adskilligt
lettere ved at føle sig helt løsrevet fra Civilisationen. Da han dernæst
havde faaet Kamelhaarsskjorten paa, lod han sig ikke fotografere i den.

\section*{III}

Paa sit Højdepunkt angriber Tolstoj da nu sin Samtids Mennesker i Kraft
af Forhold, der har deres Grund i at Mennesket er et til Menneskelighed
udviklet Dyr. Kunde han, vilde han gerne helt udrydde Dyret i Mennesket.

Det er forunderligt med denne Betegnelse: Dyret. Det er jo et ældgammelt
kristeligt Udtryk. Man glemmer ikke let den levende Skildring af Dyret i
Johannes Aabenbarings trettende Kapitel: «Jeg stod paa Havets Bred og saa
et Dyr stige op af Havet. Det havde ti Horn og syv Hoveder; paa dets Horn
var ti Diademer og paa dets Hoveder bespottelige Navne, og Dyret, jeg
saa, var ligt en Leopard; dets Fødder var som en Bjørns, og dets Svælg
som en Løves Svælg, og Dragen gav det sin Kraft og de tilbad Dyret og
sagde: hvem er som Dyret og hvem kan kæmpe mod det?»

Dyret var for Datidens Kristne som bekendt nærmest den romerske
Verdensmagt; vi véd det med al Nøjagtighed; Sindbilledet ligger
gennemsigtigt for os lige indtil det hemmelighedsfulde Tal 666 i Vers 18,
som er Neros Navn, Nero Kesar, skrevet med hebraiske Bogstaver og
sammenlagt efter disse Bogstavers Talværdi. Men dybere opfattet var Dyret
selvfølgelig for de første Kristne Hedenskabet som Væsen og Magt, alt det
Udøbte, forsaavidt Dyriske i Mennesket, som skulde overvindes. For at
bruge det moderne Kunstord: Dyret var intet virkeligt Dyr, det var
Menneskedyret.

Saaledes opfattede ogsaa Alexandre Dumas det, da han i 1870 med en
aandrig Anvendelse af det gamle Bibelsted fortalte om sine Forsøg paa at
undersøge Mennesket i den store Smeltedigel Paris. Han hævdede, at
Mennesket, endog Pariseren, altid indeholdt en nok saa lille Brøkdel af
Sjæl, omtrent som den tresindstyvende homøopatiske Opløsning endnu
indeholder et Atom af den oprindelige Vædske. Og han fortalte, hvorledes
han saa mandlige og kvindelige Fæhoveder dannes og formes af Diglen
Dampe, da han med Et hørte en kogende Lyd, og op af Kedelen steg, ikke
formet af dens Skum eller Em, men dannet af selve Stofferne i den, et
uhyre Dyr med syv Hoveder og ti Horn, og paa disse Horn ti Diademer og
paa disse Hoveder Haar med Metalglans og Alkoholfarve. Dyret lignede en
Leopard, dets Fødder var stærke som en Bjørns, dets Svælg var som en
Løves Svælg og Dragen gav det sin Styrke. Og Dyret var klædt i Purpur og
Skarlagen og pyntet med Guld, med Edelstene og Perler, og i sine hvide
Hænder holdt det, som man bærer en Skaal med Mælk, en Guldvase fuld af
alle Babels, Sodomas og Lesbos's Afskyeligheder og Urenligheder. Fra dets
Legeme udgik en berusende Damp, gennem hvilken det skinnede som den
skønneste af Guds Engle, og i hvilken tusind smaa Mennesker bevægede sig,
vred sig af Vellyst og hylede af Smerte og forsvandt med et lille Smeld
eller Knald, det vil sige: de revnede og intet blev tilbage af dem uden
en Draabe af noget Flydende, en Taare eller en Blodsdraabe. Men Dyret
blev ikke mæt. Det knuste dem med sine Fødder, sønderrev dem med sine
Negle, tyggede dem med sine Tænder, kvalte dem mod sit Bryst, og de,
hvem det saaledes kvalte, var de mest misundte. Dets syv Hoveder dannede
en Krans, der naaede Himlen; dets syv Munde var stedse smilende og
Læberne luerøde, og ovenover dets ti Diademer flammede i en Lysglans det
ene Ord: Prostitution.

Dette Dyr, det var, siger han, Kvindens nyeste Legemsform. Dyret er da
her det efter Mændene smilende og brølende Køn, der efter tusindaarig
Trældom og Afmagt hævner sig paa Manden, væbnet med sin Skønhed og sine
Skønhedsmidler, og skænker ham den Kærlighed, der nedbryder og ødelægger
ham. Og det er dette Dyr, som Dumas kastede op paa Scenens Brædder, da
han rejste det store, grove Sindbillede Césarine i \textit{La femme de
Claude}. For Dumas er det da omtrent som for Tolstoj Kulturen, nærmere
Overkulturen, der gør Mennesket i bibelsk Forstand dyrisk, det vil sige
lastefuldt; og Mennesket var det efter Forudsætningen ikke, saa længe det
ikke var moderne Menneske, men underkastet overjordiske Myndigheders
ubetingede Tugt.

\section*{IV}

Samtidigt med at Tolstoj i Rusland udarbejdede \textit{Kreutzersonaten},
skrev Zola i Frankrig sin sidste Roman \textit{Menneskedyret}. Her er
Forholdet set fra en helt forskellig Side. For Zola er Dyret ikke noget,
som Civilisationen udvikler i os, men noget, der er blevet tilbage i os
som slumrende Levning fra Urtiden og som under visse Omstændigheder til
vor Rædsel kan vaagne.

Zolas Standpunkt er ikke det moralske, men det naturhistoriske. Han
fremstiller Forholdet, som det stiller sig for ham, uden Anklage, men med
en storstilet Alvor.

I Grunden er det altid Menneskedyret, der har sysselsat Zola. Visse
yngre Skribenter som Bourget og Rod interesserer sig kun for
Menneskevæsenet i dets sidste og højeste Udfoldelse under en forfinet
Civilisation: Zola griber altid tilbage til de oprindelige Instinkter og
Sansninger, af hvilke Urmennesket bestaar. Man tænke eksempelvis paa hans
Skildring af Bondekærligheden til Jorden, til Ejendom, til Vinding i
\textit{La terre} eller af Selvopholdelsesdriften i dens naiveste Form i
\textit{Le ventre de Paris}.

Her har han grebet tilbage til et af de uhyggeligste, oprindelige Hang:
Tilintetgørelsesdriften, det grusomme Instinkt, der fører til Ødelæggelse
og Mord, og har opfattet dette Hang i dets hemmelighedsfulde Sammenhæng
med den erotiske Drift. Medens Tolstoj i sin Novelle hovedsageligt
studerer Kønsttilbøjeligheden i dens Overudvikling og viser, hvorledes
Mordgalskaben spirer ud af den, idet Tilintetgørelsesdriften opstaar som
et blot Udslag af Skinsygen, fordyber Zola sig hovedsageligt i
Mord-Instinktet og skildrer en hel Række Mord og Mordforsøg, men
efterviser i anden Linje i dem alle Sammenhængen med en lavstammet, men
livskraftig Erotik.

Vi ser den oprindelige Tilbøjelighed til Drab levende hos alle de
Mennesker, han fremfører. En Ægtemand, som erfarer, at en gammel
Vellystning har bemægtiget sig hans Hustru, da hun var ung Pige, hævner
sig efter et Anfald af tilbageskuende Skinsyge ved et Mord udført i en
Jernbanevogn og tvinger sin Hustru til at hjælpe sig ved Mordet.
Bloddaaden hviler trykkende over Ægtefællerne og skiller dem snart helt
ad.

Den unge Kone tager en Elsker og fordrer at han skal dræbe hendes Mand.
Han samtykker bevidst, men i selve det Øjeblik, da han oppebier
Ægtemandens Komme for at myrde ham, dræber han, ledet af en
uimodstaaelig Drift, imod sin Vilje sin Elskerinde.

Denne unge Mand, den ufrivillige Morder, er Bogens Hovedperson. Zola
har, som han udtrykkelig har udtalt, i Modsætning til Russerne Dostojevski,
Tolstoj villet vise, at det ikke er med Forsæt og klar Bevidsthed, men
af Drift at man myrder. Han har villet levere en Art Modstykke til
Raskolnikov, og træffer nu mærkværdigt sammen med Tolstoj.

Jacques Lantier lider af en sygelig Drift, et hemmelighedsfuldt og
uimodstaaeligt Hang, som findes i Naturen, er Lægerne vel bekendt og
hyppigt er beskrevet i lærde Haandbøger. Men hvad Zola ud af sit Eget
har tilføjet som personlig Særhed, som usandsynlig Forklaring, det er
dette: Hvad der saaledes kommer over Lantier er ikke simpelthen en
pludselig Betagelse af blindt Raseri, men en stadigt genkommende Tørst
efter at hævne meget gamle Krænkelser, om hvilke han har mistet al
nøjagtig Erindring. Det kom saa langt borte fra, hedder det, fra alt det
Onde, Kvinder havde gjort hans Forfædre, fra det Nag, der havde samlet
sig op fra den ene Han til den anden, lige fra en Kvinde bedrog en af dem
første Gang, medens de endnu beboede Huler. Anfaldene betragtes da
aldeles som Gengangervæsen; de stammer fra den forhistoriske Tid, ganske
modsat Posdnysjevs Anfald hos Tolstoj, der udledes af de moderne Tiders
Overkultur. Naar Lantier har sine Anfald, føler han det som en
Nødvendighed at kæmpe for at erobre en Kvinde og have hende helt for sig.
Dette er for Zolas Fantasi det vanslægtede Hang til at kaste hende over
sin Ryg som et Bytte, Manden for bestandig har fravristet Fællesskabet
med de andre.

Hos Zola altsaa slet fordøjet Darwinisme som hos Tolstoj fornuftstridigt
genfødt Oldkristendom.

Om Lantier er nu grupperet andre Personligheder med lignende Muligheder
til Mordgerning: En ung Pige, en Banevogterske, der af Skinsyge lader et
helt Jernbanetog smadres, fordi hun er forelsket i Lantier og vil hævne
sig paa ham og hans Elskerinde, der begge befinder sig i Toget; en gammel
Mand, et moralsk Uhyre, der langsomt forgiver sin Kone for at faa Tag i
hendes Sparepenge; endelig en Fyrbøder, der af Skinsyge tilsidst myrder
Jacques.

Dette er da et stort Digt om forhistorisk Mordlyst i dens Forbindelse med
Elskovsdriften; forsaavidt en Art forhistorisk Epos. Det er lærerigt at
sammenligne den med et virkelig forhistorisk Epos fra vore Dage som
f.\ Eks.\ Tyskeren Heinrich Harts \textit{Tul og Nahila}, det første Bind
af et Værk \textit{Lied der Menschheit}, som efter den naive Plan skal
bestaa af ikke mindre end 24 Sange.

Harts Digt foregaar paa Ceylon i den graa Oldtid Aartusinder før al
Civilisation. Hovedpersonerne er et Menneskepar, der lever iblandt
Væsener fra Overgangstiden mellem Menneskenes Forfædre og Menneskene
selv. Abeagtige Dværge og Kæmper driver endnu deres uhyggelige Spil ved
Siden af Menneskene. Kun faa Begreber har udviklet sig: Ilden er Slange,
Ildslangen; hvorledes man gør sig til Herre over den, vides endnu ikke.
Man bruger Udtryk som «Ildslangen har naget paa Kødet». Faa Redskaber:
Kølle, Stenøkse, Stenspids til Spydet, Fiskekrog af Ben, Bastslynge til
Fangst. Blandt Menneskene raader Naturafhængighed og Hjordfølelse. Vi
følger Menneskeparrets Livshistorie. Nahila elsker Tul for Livet, fordi
denne Mand for hendes Skyld blev udstødt af sin Stamme, idet han ikke
vilde prisgive hende til Fællesskabet. Skamfølelse, Undseelse udvikles
hos hende af en ejendommelig og ikke usandsynlig Bevæggrund: fordi hendes
eget Køn synes hende hæsligt: derfor laver hun sig et Skørt af Blade. Vi
ser, hvorledes Barnets Fødsel under bestemte Forhold gør Husdyrtugt
nødvendig. En Ged indfanges, der nærer Barnet med sin Mælk. Saa store
Vanskeligheder Emnet end frembyder, er Digtet ikke dumt. Menneskene synes
her virkeligt oprindelige og vilde. Og dog er Alt hos Zola i sin Art
vildere og større, skønt det foregaar i vore Dage, i 1870, og fremtræder
som en Roman om Jernbanedrift og Jernbane-Uheld og om Forbrydelser, der
begaas med Jernbaneskinner eller i Jernbanevogne.

Jules Lemaître har gjort opmærksom paa den Svaghed i denne Bog, at
Omgivelserne ikke har noget egenligt med Kærnen at gøre. Medens der i
\textit{L'Assommoir} er en nødvendig Sammenhæng mellem Brændevinskippen
og Arbejderbefolkningen eller i \textit{Germinal} mellem Gruben og
Minearbejderne, finder der her ingen nødvendig Sammenhæng Sted mellem
Tilintetgørelseslyst og Jernbanedrift. Forsaavidt er denne Bog ringere
end adskillige foregaaende, mere tilfældig i sin Komposition. Men den er
stor ved Digterens Blik for de simple og skrækkelige Grunddrifter i
Menneskeheden og stor ved sin rige Sindbilledkunst.

Foruden Menneskedyret er der i denne Roman som i Zolas andre et
upersonligt Væsen, der tilegner sig en Hovedlod i Forfatterens Interesse.
I tidligere Bøger forekom Kirkegaarden, Hallerne, Brændevinskænken,
Salgsmagasinet, Agerjorden, Kirken. Vi har ovenfor set, at medens Zola
gennemgaaende fører det Menneskelige tilbage til det Dyriske, tillægger
han upersonlige Skabninger al den mere end dyriske Kraft, han berøver
Enkeltværenerne, ja Selvstændighed og Vilje. En saadan upersonlig
Skabning er i Jernbaneromanen et Lokomotiv. Det har et Navn; det hedder
Lison, fører et Liv og vækker Følelser som en Kvinde. Lokomotivføreren
elsker sin Maskine, elsker den som Rytteren sin Hest. Der finder en Art
Ægteskab Sted imellem dem.

Her forekommer en storladen Skildring af Lokomotivets Kamp for at trænge
igennem ophobet Sne under Uvejr, og da Toget senere hen støder mod en
Karre, fuld af svære Sten, og Lokomotivet sprænges, saa dør det ganske
som et levende Væsen. Der er heri en Kunst, som er bleven til Kunstleri.
Lokomotivet, dette Uhyre med gloende Øjne og Jernlænder, gør saa at sige
Menneskene til Skamme ved sin Lærvillighed, sin Mandstugt, sin hele
Skikkelighed ligesom ved den ublandede Kærlighed, det indgyder sin Fører.

Og ved Bogens Slutning savnes ikke heller et andet, større Sindbilled. Da
Fyrbøderen har myrdet Jacques, stødt ham ned fra hans Lokomotiv, og da
dette uden Fører, uden ledende Aand som et rasende Væsen farer Station
efter Station forbi, førende franske indkaldte Soldater til Grænsen i
Juli Maaned af hint mærkværdige Aar 1870, saa staar der:

\begin{quote}
«Det rullede, rullede afsted i den sorte Nat, Ingen vidste hvorhen.
Hvad gjorde det, hvor mange Ofre Maskinen knuste paa Vejen! Gik den ikke
trods Alt mod Fremtiden, ubekymret om det udgydte Blod? Uden Fører, midt
i Mørket, som et blindt og døvt Dyr, rullede, rullede den, dragende efter
sig al denne Kanonføde, disse allerede udmattede Soldater, der i deres
Drukkenskab sang.»
\end{quote}

Dette Lokomotiv, som ikke ledes af Nogen, det er, som man føler,
Sindbilledet paa hint Frankrig, der drager i Krig uden en Vilje i sin
Spidse. Romanen om Tilintetgørelsesdriften, om Mord-Instinktet, naaer sit
Højdepunkt i denne Henslæben af Tusinder og atter Tusinder til dette
uhyre Massemord. Den store Krigs endeløse Blodsudgydelse er
Menneskedyrets Triumf og Nederlag i Et.

\section*{V}

Det er det Anerkendelsesværdige ved disse store fremmede Digtere og
Skribenter, at de for vore Øjne slaar Menneskevæsenet helt ud som en
Vifte. Der er Intet, som af Frygt for Læsere eller Bladskrivere skal
borthykles eller bortlyves. Vi ser ud over et vidt Felt af sjæleligt Liv
fra det Lastefulde og Lave, fra Forbrydelsen, over det dyrisk
Menneskelige, det blot Naturlige, det Altformenneskelige til det rent
Menneskelige og dets Menneskeadel. Og alt, hvad der meddeles, er meddelt
i de store Literaturers beundringsværdige Frisprog uden Pænhed og uden
Følsomhed.

Alle behøver de for at give en Forestilling om Livets Rigdom og
Elendighed Forestillingen om en Spaltning af vort Væsen i to Verdner, den
ubevidste og den bevidste, der i den seneste Tids poetiske Sprog er blevne
gjorte enstydige med Dyrets og Menneskets Verdener i os,
Nødvendighedens og Frihedens Omraader, som man i gamle Dage sagde.

Den moderne Videnskab tror ikke mere paa en Viljesfrihed som den, man
tidligere klamrede sig til. Enkelte af Digterne ser netop heri en Fare.
Saaledes Bourget, der synes skræmt af det Indtryk den filosofiske
Nødvendighedslære undertiden har gjort paa en Ungdom, i hvis Øjne den
giver et Fribrev til hensynsløs Egenkærlighed. Han frygter, at selve
Videnskaben, selve Civilisationen skal slippe Dyret løs i de kommende
Slægter og han har i sine sidste Bøger (\textit{Mensonges}, \textit{Le
Disciple}) ikke uden en vis fromladen Slaphed henvist til den katolske
Abbed, til Religionen, til Bønnen, som til den eneste Magt, der kan holde
Dyret nede og give de moderne Sind Enhed paany.

Dumas, der er ægte Franskmand nok til ikke at nære nogen Tvivl om Viljens
Frihed, ser Dyret udvikle sig i Overcivilisationens Morads og raaber Vagt
i Gevær imod det. Moralisten Tolstoj og Verdensbarnet Maupassant ser med
samme Lede paa det Grundforhold imellem de to Køn, der for dem begge
staar som en Krænkelse af Menneskeværdigheden, men der for Maupassant er
nedværdigende som et Stykke uovervunden og uovervindelig Urnatur, for
Tolstoj nedværdigende som et Foster af Kulturens overophidsede Unatur.
Zola endelig er den gennemført overbeviste Nødvendighedstroende, for hvem
Menneskeplanten, Menneskedyret, Mennesketanken, Menneskeviljen lyder
uafrystelige Love, og som til Gengæld finder en Art fantastisk
Tilfredsstillelse i at udruste livløse Genstande med selvstændigt Liv og
uafhængig Vilje.

Menneskelig, virkeligt menneskelig er kun Maskinen, synes Zola at sige i
sin sidste Roman. Menneskelig, virkeligt menneskelig er kun Civilisationen,
synes Maupassant at sige i sin sidste Novellesamling, der bryder Staven
over alt det, der er Naturen i os, som dyrisk. Menneskelig, adeligt
menneskelig, i gamle Dages Sprog guddommelig, er kun Sjælens Samling i
from Følelse og Bøn, siger Bourget i sine to sidste Bøger. Menneskelig er
ikke Kulturen, ikke Civilisationen, siger Tolstoj i sine sidste Værker,
men kun den Levevis, der aldrig afbøder Ondt med Nødværge, aldrig
afværger Uret, aldrig bekymrer sig om Ejendom og Penge, og som bringer
Mand og Kvinde til altid i hinanden kun at se en Søster og Broder.

Det er da det Sprog, som føres af vore Dages Vise, af de Mænd, der
nutildags henvender sig til den største og ypperste Læsekres paa Jorden.

Hvis man kunde tænke sig, at en gammel Hellener fra Grækenlands klassiske
Tid opstod af sin Aske og hørte disse Udsagn, vilde han sikkert slaa
Hænderne sammen af Forbavselse over saadan Tale.

\end{document}
